\documentclass[12pt,a4paper]{article}
\usepackage[margin=1in]{geometry}
\usepackage{amsmath,amssymb}
\usepackage{graphicx}
\usepackage{booktabs}
\usepackage{array}
\usepackage{float}
\usepackage{hyperref}
\usepackage{enumitem}
\usepackage{caption}
\usepackage[T1]{fontenc}
\usepackage[utf8]{inputenc}

\title{\textbf{Deep Learning Based CT Reconstruction Enhancement}\\
Technical Communication Project Report}
\author{
Hrushikesh Waghmare (23324025) \and
Pushkar Singh (23411026)
}
\date{\today}

\begin{document}
\maketitle

\begin{abstract}
This report presents the work completed for the Technical Communication project on CT image reconstruction and image enhancement. The project starts from raw X-ray projection data, reconstructs a 3D CT volume using classical methods, creates degraded acquisition scenarios such as sparse-view, limited-angle, and noisy scans, and prepares paired training data for a U-Net based enhancement model. The core contribution of the work is a complete pipeline that combines physics-based reconstruction with supervised deep learning for artifact reduction and quality improvement.
\end{abstract}

\section{Introduction}
Computed tomography (CT) reconstructs the internal structure of an object from multiple X-ray projections acquired at different angles. In ideal conditions, a large number of clean projections are available, which leads to a high-quality reconstruction. In practical scenarios, however, the scan may suffer from fewer views, limited angular coverage, or noise. These limitations produce streaking, blurring, and missing-structure artifacts in the reconstructed image.

The aim of this project is to address this problem using a hybrid approach:
\begin{enumerate}[leftmargin=*]
\item reconstruct a reference CT volume from full projection data,
\item simulate degraded acquisition conditions,
\item reconstruct artifact-corrupted volumes from these degraded projections,
\item train a U-Net model to enhance the degraded reconstructions.
\end{enumerate}

The scope of this report is intentionally limited to CT reconstruction and U-Net based image enhancement. Auxiliary export and STL/phantom-generation utilities are excluded because they are outside the central scope of the Technical Communication submission.

\section{Problem Formulation}
Let $f(x,y,z)$ denote the unknown 3D attenuation distribution of the scanned object. The CT scanner measures projection data $g$ obtained through the forward projection operator $\mathcal{P}$:
\[
g = \mathcal{P}f + \epsilon,
\]
where $\epsilon$ represents measurement noise and acquisition imperfections.

The reconstruction task is to estimate the object from the measured projections:
\[
\hat{f} = \mathcal{R}(g),
\]
where $\mathcal{R}$ is a reconstruction algorithm.

For degraded acquisition, the measured data are modified by a degradation operator $\mathcal{D}$:
\[
g_d = \mathcal{D}(g).
\]
The degraded reconstruction is therefore
\[
\hat{f}_d = \mathcal{R}(g_d).
\]

The deep-learning objective is to train a model $\mathcal{N}$ such that
\[
\mathcal{N}(\hat{f}_d) \approx \hat{f}_{\text{ref}},
\]
where $\hat{f}_{\text{ref}}$ is the reconstruction obtained from the full projection dataset.

\section{Dataset and Geometry Handling}
The raw data are loaded from the folder \texttt{sample 1}. The script \texttt{data\_loader.py} reads:
\begin{itemize}[leftmargin=*]
\item all TIFF projection images,
\item the scanner configuration file \texttt{settings.cto}.
\end{itemize}

The projection stack is stored as a 3D array of shape
\[
(\text{number of projections}, \text{detector rows}, \text{detector columns}).
\]

The geometry parser extracts the main CT parameters:
\begin{itemize}[leftmargin=*]
\item number of projections,
\item angular range,
\item detector pixel size,
\item source-to-object distance (SOD),
\item source-to-detector distance (SDD),
\item center of rotation,
\item valid axial reconstruction interval.
\end{itemize}

The projection angles are generated as
\[
\theta_i = \theta_0 + i\Delta\theta,\qquad i=0,1,\dots,N-1,
\]
with
\[
\Delta\theta = \frac{\text{angular range}}{N}.
\]

\section{Classical Reconstruction Approach}
\subsection{Slice-wise FBP Baseline}
The first reconstruction approach implemented in the project is a slice-wise filtered backprojection (FBP) baseline. This method was useful for validating the pipeline and understanding the relationship between projections, sinograms, and reconstructed slices.

For a chosen detector row, one line is extracted from every projection to form a sinogram:
\[
S(s,\theta),
\]
where $s$ is detector position and $\theta$ is the projection angle.

The sinogram is filtered in the frequency domain using the ramp filter:
\[
\tilde{S}(s,\theta)=\mathcal{F}^{-1}\left(|\omega|\mathcal{F}(S(s,\theta))\right).
\]

The reconstructed slice is then obtained by backprojection:
\[
f(x,y)\approx \int \tilde{S}(x\cos\theta + y\sin\theta,\theta)\, d\theta.
\]

Although this baseline is not the final solution for cone-beam CT, it confirms that the data loading, angle generation, and reconstruction logic are working correctly.

\begin{figure}[H]
\centering
\includegraphics[width=0.8\textwidth]{outputs/fbp_baseline/fbp_preview.png}
\caption{Preview generated by the slice-wise FBP baseline. The figure shows axial, coronal, and sagittal views of the reconstructed volume subset.}
\end{figure}

\subsection{Main 3D Cone-Beam FDK Reconstruction}
The main reconstruction method used in this project is a 3D cone-beam Feldkamp-Davis-Kress (FDK) reconstruction implemented with the ASTRA toolbox on CUDA.

Before reconstruction, the raw detector data are spatially downsampled to reduce GPU memory usage. If the downsampling factor is $k$, block averaging is used:
\[
P_{\text{ds}}(i,j)=\frac{1}{k^2}\sum_{u=0}^{k-1}\sum_{v=0}^{k-1} P(ki+u,kj+v).
\]

The projection intensities are converted into attenuation values using the Beer-Lambert relation:
\[
\mu L = -\log\left(\frac{I}{I_0}\right),
\]
where $I$ is the measured intensity and $I_0$ is the air intensity estimate.

The voxel size is estimated from scanner geometry as
\[
\text{voxel size} = \text{detector pixel size}\cdot \frac{\text{SOD}}{\text{SDD}}.
\]

For the saved full reconstruction in this project, the metadata show:
\begin{itemize}[leftmargin=*]
\item 359 input projections,
\item downsampled projection shape: $(359, 500, 500)$,
\item reconstructed output volume shape: $(441, 500, 500)$,
\item voxel size: approximately $0.0768$ mm.
\end{itemize}

\begin{figure}[H]
\centering
\includegraphics[width=0.8\textwidth]{outputs/fdk_astra/fdk_preview.png}
\caption{Main cone-beam FDK reconstruction preview generated using ASTRA.}
\end{figure}

\section{Simulation of Degraded CT Acquisition}
To train a deep model for image enhancement, degraded acquisition scenarios were simulated from the original projection dataset. This provides paired examples of low-quality and reference reconstructions.

\subsection{Sparse-view CT}
Sparse-view CT reduces the number of projection angles by keeping only every $k$-th projection:
\[
g_{\text{sparse}} = \{g_0, g_k, g_{2k}, \dots\}.
\]
This leads to severe streak artifacts due to under-sampling.

In the prepared datasets:
\begin{itemize}[leftmargin=*]
\item \texttt{sparse\_view\_step\_2} keeps 180 projections,
\item \texttt{sparse\_view\_step\_4} keeps 90 projections.
\end{itemize}

\subsection{Limited-angle CT}
Limited-angle CT keeps only projections in a restricted angular interval:
\[
\theta \in [\theta_{\text{start}},\theta_{\text{stop}}].
\]
This creates incomplete data and directional artifacts.

In the prepared datasets:
\begin{itemize}[leftmargin=*]
\item \texttt{limited\_angle\_0\_180} keeps 180 projections,
\item \texttt{limited\_angle\_30\_210} keeps 180 projections.
\end{itemize}

\subsection{Noisy CT}
Noise was also added to the projection data. The main saved datasets use Poisson noise, which is appropriate for X-ray photon-counting processes. In simplified form:
\[
c \sim \text{Poisson}(\lambda),
\]
followed by scaling back to intensity space.

The prepared noise datasets are:
\begin{itemize}[leftmargin=*]
\item \texttt{poisson\_noise\_0.25},
\item \texttt{poisson\_noise\_0.5}.
\end{itemize}

\begin{figure}[H]
\centering
\begin{minipage}[t]{0.48\textwidth}
\centering
\includegraphics[width=\textwidth]{outputs/degradations/sparse_view_step_4/preview.png}
\captionof{figure}{Example degraded projection data preview for sparse-view CT with step 4.}
\end{minipage}
\hfill
\begin{minipage}[t]{0.48\textwidth}
\centering
\includegraphics[width=\textwidth]{outputs/degradations/limited_angle_0_180/preview.png}
\captionof{figure}{Example degraded projection data preview for limited-angle CT from $0^\circ$ to $180^\circ$.}
\end{minipage}
\end{figure}

\section{Training Pair Generation}
The script \texttt{build\_training\_pairs.py} reconstructs both the full dataset and the degraded datasets using the same FDK pipeline. This creates supervised pairs of the form:
\[
\text{degraded reconstruction} \rightarrow \text{reference reconstruction}.
\]

If the degraded and reference volumes differ slightly in shape, they are aligned by center-cropping to a shared size:
\[
(Z,Y,X)_{\text{shared}} = \min\left((Z,Y,X)_{\text{input}}, (Z,Y,X)_{\text{target}}\right).
\]

Each volume is then normalized to the interval $[0,1]$:
\[
v_{\text{norm}}=\frac{v-v_{\min}}{v_{\max}-v_{\min}}.
\]

Finally, axial slices are saved as paired arrays for deep-learning training.

From the saved metadata, six aligned training-pair volumes were prepared:
\begin{itemize}[leftmargin=*]
\item sparse-view step 2,
\item sparse-view step 4,
\item limited-angle $0^\circ$ to $180^\circ$,
\item limited-angle $30^\circ$ to $210^\circ$,
\item Poisson noise 0.25,
\item Poisson noise 0.5.
\end{itemize}

Each saved pair has aligned volume shape
\[
(441, 500, 500),
\]
which means the prepared dataset corresponds to
\[
6 \times 441 = 2646
\]
axial slice pairs in total.

\begin{figure}[H]
\centering
\begin{minipage}[t]{0.48\textwidth}
\centering
\includegraphics[width=\textwidth]{outputs/training_pairs/sparse_view_step_4/input_fdk_preview.png}
\captionof{figure}{Input reconstruction from the degraded sparse-view step 4 dataset.}
\end{minipage}
\hfill
\begin{minipage}[t]{0.48\textwidth}
\centering
\includegraphics[width=\textwidth]{outputs/training_pairs/sparse_view_step_4/target_fdk_preview.png}
\captionof{figure}{Reference target reconstruction from the full projection dataset.}
\end{minipage}
\end{figure}

\section{Deep Learning Model for Image Enhancement}
The enhancement model implemented in \texttt{train\_unet.py} is a 2D U-Net trained on axial slice pairs.

\subsection{Why U-Net}
U-Net is suitable for image restoration because it combines:
\begin{itemize}[leftmargin=*]
\item an encoder path that captures context,
\item a decoder path that reconstructs details,
\item skip connections that preserve fine spatial information.
\end{itemize}

Each input slice and target slice has the form
\[
x,y \in \mathbb{R}^{1\times H\times W}.
\]

The implemented network uses convolutional blocks and an encoder-decoder structure with skip connections. The default channel sizes are:
\[
(32, 64, 128, 256).
\]

\subsection{Loss Function}
The model is trained using L1 loss:
\[
\mathcal{L}_{L1} = \frac{1}{N}\sum_{i=1}^{N}\left|\hat{y}_i-y_i\right|.
\]
This loss is a reasonable choice for reconstruction enhancement because it penalizes absolute deviations while preserving edges better than many overly smooth alternatives.

\subsection{Validation Metric}
The script also computes peak signal-to-noise ratio (PSNR):
\[
\text{PSNR}=10\log_{10}\left(\frac{1}{\text{MSE}}\right),
\]
for normalized image values in the range $[0,1]$.

\subsection{Training Configuration}
The implemented default training configuration is:
\begin{itemize}[leftmargin=*]
\item epochs: 20,
\item batch size: 8,
\item learning rate: $10^{-3}$,
\item validation fraction: 0.2.
\end{itemize}

Therefore, the completed work is not only data preparation but also a full training-ready enhancement pipeline for degraded CT reconstructions.

\section{Overall Workflow}
The complete workflow of the project can be summarized as follows:
\begin{enumerate}[leftmargin=*]
\item load raw projection images and CT geometry,
\item reconstruct a baseline volume using slice-wise FBP,
\item reconstruct the full reference volume using 3D FDK,
\item simulate sparse-view, limited-angle, and noisy acquisition conditions,
\item reconstruct degraded volumes from the simulated projection sets,
\item align degraded and reference volumes,
\item extract normalized axial slice pairs,
\item train a 2D U-Net for image enhancement.
\end{enumerate}

\section{Conclusion}
This project demonstrates a complete and technically coherent pipeline for CT image enhancement. The work begins with raw X-ray projections and scanner geometry, proceeds through classical CT reconstruction, generates realistic degraded acquisition settings, and prepares supervised training data for a U-Net based enhancement model.

The main achievement of the work is the integration of physics-based reconstruction and deep learning. The reconstruction stage provides interpretable reference and degraded volumes, while the U-Net stage provides a practical method for improving image quality under difficult acquisition conditions. This makes the project relevant both from an imaging perspective and from a machine-learning perspective.

\section{Future Scope}
The most relevant future extensions of this work are:
\begin{itemize}[leftmargin=*]
\item training completion and quantitative evaluation on saved checkpoints,
\item adding SSIM and error-map analysis in addition to PSNR,
\item performing full-volume inference after training,
\item comparing U-Net enhancement with iterative reconstruction methods,
\item extending from 2D slice-wise enhancement to 3D deep-learning models.
\end{itemize}

\end{document}
